\bookchapter{Chapter 9: Mobile Office and Learning}

\chapterintro{Microsoft 365}

\booksection{: If your company uses genuine Office, this is the best choice. Format compatibility with the desktop version is nearly perfect, avoiding the slight layout shifts that WPS may have with some complex layouts. Deeply integrated with OneDrive, the document you were halfway through editing on your computer opens as the latest version on your phone.}

Shimo Docs / Tencent Docs

\begin{itemize}
  \item \textbf{: These two represent "online collaboration". They are suitable for multiple people editing the same document simultaneously, with each person's changes shown in real time. I used Shimo Docs to plan a trip with a friend — both of us filling in information into a spreadsheet at the same time was far more efficient than doing it alone.}Practical Workflow
  \item \textbf{: My current habit is — important files go to OneDrive, daily editing uses WPS, and team collaboration opens Shimo Docs. Since my files went to the cloud, I've never again had the embarrassment of "forgetting my USB drive".}2. Note-Taking Apps: From Quick Notes to a Knowledge System
  \item \textbf{When it comes to taking notes, many people want to build a complex system right from the start, only to give up after three days. My advice is to start light and then go heavy — first build the habit of "recording", then consider the problem of "managing".}Quick Notes — The built-in system is enough
\end{itemize}

\textbf{: iOS Notes launches extremely fast and supports document scanning, handwriting, and audio recording. Android's built-in memo apps by various brands are also decent, and can be launched quickly even from the lock screen. These tools are suitable for recording temporary ideas, shopping lists, and meeting key points. The key is fast launch — you can record as soon as you think of something, without waiting for an app to load.}Read It Later — Collecting Information

\booksection{: WeChat's “Float” feature is actually a hidden gem. When you see a good article but don't have time to read it carefully, tap “Float” in the upper-right corner, and it shrinks into a small dot attached to the edge of the screen, ready to be tapped open to continue reading anytime. It's much more convenient than forwarding to “File Transfer Assistant”. In addition, read-later tools like Cubox or Pocket can collect good content you come across anywhere, across platforms.}

Building a System — Notion and Feishu

\begin{itemize}
  \item \textbf{: When you feel that scattered notes are becoming hard to find, it's time to move to Notion. This tool is truly flexible — you can build a "reading notes database" that automatically organizes by book title, author, and rating; or set up a "project board" to visualize the progress of work tasks. Feishu's "Bitable" is similar, and has faster access speeds in China. I use Notion to manage all my writing materials, study notes, and project progress — I've been using it for over three years.}A Practical Note-Taking Strategy
  \item \textbf{: My own habit is a "three-piece set" — temporary ideas go to the system memo (instant launch), articles to read go to WeChat's float (handy), and knowledge that needs long-term accumulation goes to Notion (systematized). Don't try to stuff everything into one tool from the start — that makes it easier to give up.}3. Scanning and OCR: Bringing the Paper World into Your Phone
  \item \textbf{Phone scanning, when used well, can really save a lot of trips to the print shop. Last time I applied for a visa and needed electronic copies of a bunch of materials, I got them all done with my phone.}iPhone users
\end{itemize}

\textbf{: The built-in scanning feature of iOS Notes is already quite good. Open Notes tap the camera icon select “Scan Documents”, aim at the document and it will automatically recognize the border and correct the perspective. After shooting, you can adjust the color — the “Black \& White” mode produces results closest to a scanner, and “Grayscale” is suitable for materials with images.}Android users

\booksection{: Many brands' cameras have a built-in “Document Scan” mode. If you can't find it, download Microsoft Lens from the app store (free, and integrated with OneDrive).}

OCR Text Recognition

\begin{itemize}
  \item \textbf{: If you need to extract text from scanned documents for editing, you need OCR. WeChat's built-in "Extract Text" feature is the simplest — long-press the image to use it, and the recognition accuracy is decent. For bulk processing, you can use the Baimiao app or CamScanner. For Chinese recognition, Baidu's OCR engine works quite well; WPS's image-to-text feature is also very accurate.}Scanning Tips
  \item \textbf{Look Up Tips}4. How to Choose a Cloud Drive: There's No Best, Only What's Best for You
  \item \textbf{I thought about choosing a cloud drive for a long time, and finally realized it really depends on what your "primary ecosystem" is.}Apple Ecosystem iCloud
  \item \textbf{: Seamless syncing that makes you barely notice its existence. Photos you take on your phone automatically appear on your computer a few minutes later. And Apple's privacy protection is the strictest. But the drawbacks are also obvious: cross-platform sharing with Android or Windows users is cumbersome, the free tier is only 5GB, and expansion costs money.}Windows Primary OneDrive
\end{itemize}

\booksection{: If you use a Windows PC + Office 365, OneDrive is a natural match. Save a document in Word, and it's automatically stored in the cloud. The 1TB space bundled with the Office subscription is a great deal. The mobile app's auto photo backup feature is also very handy.}

Domestic Sharing Baidu Netdisk

\begin{itemize}
  \item \textbf{: To be honest, Baidu Netdisk's download speed limit has always been criticized, but its sharing ecosystem is truly unmatched. A large amount of study materials and film/TV resources on the Chinese internet are shared via Baidu Netdisk. If you only occasionally download files shared by others, you can make do without a membership. If you're a heavy resource collector, just pay for the membership — the time saved is worth much more.}Aliyun Drive / Quark Drive
  \item \textbf{: These are the rising stars in recent years. Aliyun Drive has no upload/download speed limits (currently), and Quark Drive is deeply integrated with the Quark browser, so you can directly save to the drive when searching for materials, which is very convenient.}My Combined Solution
  \item \textbf{: iCloud manages personal photos and daily files (privacy first), Baidu Netdisk stores resources and receives shares (ecosystem first), and OneDrive manages work documents (Office integration first). Three cloud drives, each with its own job, not interfering with each other.}5. A Real Mobile Office Scenario
  \item \textbf{Let me share a real experience of my own. Once on a business trip out of town, the client suddenly needed to revise contract terms. My workflow was as follows:}At the hotel, I received the client's revision comments via WeChat (text messages + a photo of handwritten annotations).
\end{itemize}

\textbf{Opened OneDrive to find the original contract, opened and edited it with WPS, revising the terms one by one according to the client's WeChat messages.}After revisions, used WPS's “Export as PDF” feature, and sent it to the client's WeChat group for confirmation.

\booksection{The client raised a second round of comments, so I used WPS's brush tool to directly circle the revision positions on the PDF, with voice annotations explaining the revision approach.}

After confirming everything was correct, used WPS's “Electronic Signature” feature to sign online, and sent it to the other party to sign as well.

\begin{enumerate}
  \item Finally, archived the signed contract PDF to both Baidu Netdisk and iCloud, and made a note in my work log in Notion.
  \item The whole process took about forty minutes from start to finish, without turning on a computer at all. To be honest, three years ago I would never have believed a phone could do these things, but now it really can. The key is to familiarize yourself with these tools in advance, so you don't scramble to download and register when you actually need them.
  \item Try it out: Build your mobile office workflow
  \item Choose your primary document app (WPS or Microsoft 365), log in and link it to a cloud drive.
  \item Open Notes or Memo, use the scan feature to scan a document, and export it as a PDF to send to yourself.
  \item Choose a cloud drive app and set it to auto-backup your phone photos (select "Backup only on Wi-Fi" to save data).
\end{enumerate}

Create a new page in Notion or Feishu, and try building a simple “work log” template.

\begin{tcbpractice}{Chapter 10: Mobile Payment and Life Services}
\begin{enumerate}
  \item To be honest, my parents' generation was initially resistant to the idea of "paying with a phone". "If the money is in the phone, what if it gets stolen?" This concern is very normal, but the reality is: the cash in your pocket has a much higher chance of being lost than mobile payment being fraudulently charged. Mobile payment is not about storing money in your phone — it's just a key, a key that opens your bank account.
  \item 1. How Exactly Do You Bind a Bank Card?
  \item Whether you use Alipay or WeChat, the first time you bind a card, you need to go through these steps:
  \item Open the app, find "Me" or "Wallet". Tap "Add Bank Card".
\end{enumerate}

\end{tcbpractice}
